\[ %%%%%%%%%%%%%%%%%%%%%%%%%%%%%%%%%%%%%%%%%%%%%%%%%%%%%%%%%%%%%%%%%%%%%%%%%%%%%%%% \subsection{Motivation} %%%%%%%%%%%%%%%%%%%%%%%%%%%%%%%%%%%%%%%%%%%%%%%%%%%%%%%%%%%%%%%%%%%%%%%%%%%%%%%% The rapid development of Scientific Machine Learning has fundamentally transformed the way complex physical systems are modeled and simulated. Instead of relying exclusively on computationally expensive numerical solvers, modern learning-based surrogate models increasingly aim to approximate the evolution operators of nonlinear dynamical systems directly from observational or simulation data. This paradigm has enabled substantial acceleration of scientific simulations while opening new opportunities for real-time optimization, uncertainty quantification, inverse design, and autonomous scientific discovery. However, increasing predictive accuracy alone is insufficient for many scientific and engineering applications. Numerical predictions generated by machine learning models must also remain physically meaningful throughout the entire forecasting horizon. In practical applications, including computational fluid dynamics, weather prediction, structural health monitoring, plasma simulations, digital twins, and multi-physics optimization, violations of physical consistency often become more detrimental than moderate prediction errors. Even highly accurate short-term forecasts may become unusable if accumulated numerical errors eventually produce physically impossible system states. This challenge highlights an important distinction between conventional machine learning and scientific machine learning. Traditional learning algorithms primarily seek statistical consistency with observed data distributions, whereas scientific models must additionally satisfy physical principles that remain valid independently of the available observations. Consequently, scientific learning architectures require mechanisms capable of continuously preserving physically admissible system evolution rather than merely minimizing prediction error over finite datasets. At the same time, the complexity of contemporary scientific simulations continues to increase. Modern computational models frequently involve heterogeneous spatial discretizations, irregular computational meshes, coupled multi-scale interactions, and dynamically evolving topologies. Such characteristics naturally motivate graph-based representations in which local physical interactions can be modeled directly through graph connectivity instead of regular Cartesian grids. Graph Neural Networks therefore provide an attractive computational foundation for scientific simulations because they generalize naturally across irregular domains while maintaining computational efficiency. Nevertheless, current graph-based surrogate models remain predominantly optimized from a data-centric perspective. Although graph topology effectively captures local interactions between neighboring entities, the underlying message-passing process typically remains unaware of the governing physical laws responsible for the observed dynamics. As a result, the learned latent representations may gradually drift away from physically admissible trajectories during autoregressive prediction, particularly over long temporal horizons. These limitations motivate the development of computational architectures in which physical reasoning is incorporated directly into the information propagation process rather than introduced only through external loss regularization. Such architectures should simultaneously learn latent system dynamics, preserve governing physical constraints, maintain numerical stability, and adapt to heterogeneous graph structures within a unified differentiable computational framework. \] 
